% Options for packages loaded elsewhere
\PassOptionsToPackage{unicode}{hyperref}
\PassOptionsToPackage{hyphens}{url}
\PassOptionsToPackage{dvipsnames,svgnames,x11names}{xcolor}
%
\documentclass[
  10pt,
]{article}
\usepackage{amsmath,amssymb}
\usepackage{iftex}
\ifPDFTeX
  \usepackage[T1]{fontenc}
  \usepackage[utf8]{inputenc}
  \usepackage{textcomp} % provide euro and other symbols
\else % if luatex or xetex
  \usepackage{unicode-math} % this also loads fontspec
  \defaultfontfeatures{Scale=MatchLowercase}
  \defaultfontfeatures[\rmfamily]{Ligatures=TeX,Scale=1}
\fi
\usepackage{lmodern}
\ifPDFTeX\else
  % xetex/luatex font selection
  \setmainfont[]{DejaVu Serif}
  \setsansfont[]{DejaVu Sans}
  \setmonofont[]{DejaVu Sans Mono}
  \setmathfont[]{Latin Modern Math}
\fi
% Use upquote if available, for straight quotes in verbatim environments
\IfFileExists{upquote.sty}{\usepackage{upquote}}{}
\IfFileExists{microtype.sty}{% use microtype if available
  \usepackage[]{microtype}
  \UseMicrotypeSet[protrusion]{basicmath} % disable protrusion for tt fonts
}{}
\makeatletter
\@ifundefined{KOMAClassName}{% if non-KOMA class
  \IfFileExists{parskip.sty}{%
    \usepackage{parskip}
  }{% else
    \setlength{\parindent}{0pt}
    \setlength{\parskip}{6pt plus 2pt minus 1pt}}
}{% if KOMA class
  \KOMAoptions{parskip=half}}
\makeatother
\usepackage{xcolor}
\usepackage[margin=2.0cm]{geometry}
\usepackage{longtable,booktabs,array}
\usepackage{calc} % for calculating minipage widths
% Correct order of tables after \paragraph or \subparagraph
\usepackage{etoolbox}
\makeatletter
\patchcmd\longtable{\par}{\if@noskipsec\mbox{}\fi\par}{}{}
\makeatother
% Allow footnotes in longtable head/foot
\IfFileExists{footnotehyper.sty}{\usepackage{footnotehyper}}{\usepackage{footnote}}
\makesavenoteenv{longtable}
\setlength{\emergencystretch}{3em} % prevent overfull lines
\providecommand{\tightlist}{%
  \setlength{\itemsep}{0pt}\setlength{\parskip}{0pt}}
\setcounter{secnumdepth}{-\maxdimen} % remove section numbering
% doc_header.tex -- LaTeX preamble additions used by build_docs.sh for every
% Markdown -> LaTeX -> PDF conversion (the documents of the folder Revision_HS_CJT_and_Kohn-Sham).
\usepackage{etoolbox}
% longtable may not be loaded by pandoc when a document has no table
\usepackage{longtable}
% bold vectors: \bm -> \symbf under unicode-math (XeLaTeX), bm package otherwise
\makeatletter
\@ifpackageloaded{unicode-math}{\providecommand{\bm}[1]{\symbf{#1}}}{\usepackage{bm}}
\makeatother
% smaller type in long tables so that wide audit tables fit the text width;
% inside a table a line break is allowed after an underscore of plain text
% (pandoc writes it as \_), so that long identifiers wrap inside their column
% rows are set a little apart, so that symbols with indices do not touch
\AtBeginEnvironment{longtable}{\footnotesize\renewcommand{\arraystretch}{1.45}%
  \let\sptunderscore\_\def\_{\sptunderscore\allowbreak}}
\setlength{\LTcapwidth}{\textwidth}
% A short table is kept on one page.  The build (code/scripts/keep_tables.py) ends the rows of such a table with \\*,
% which forbids a page break after a row, and sets \sptkeptpenalty to 10000 in front of the table.  This penalty is
% put here in the place of the breakpoint (penalty 0) that longtable leaves after every table, so that the paragraph
% that follows a kept table, its caption, begins on the same page; it is set back to 0 after the table.
\newcount\sptkeptpenalty
\sptkeptpenalty=0
\makeatletter
\patchcmd{\endlongtable}{\endgraf\penalty\z@\addvspace\LTpost}%
  {\endgraf\penalty\sptkeptpenalty\addvspace\LTpost\global\sptkeptpenalty=\z@}%
  {}{\PackageError{doc_header}{the end of longtable could not be patched}{The tables can still be typeset, but a short table may be separated from its caption.}}
\makeatother
% allow line breaks after slashes and hyphens inside table cells
\usepackage{xurl}
% code blocks: break over-long lines instead of running into the margin.
% Plain ``` blocks are typeset by pandoc as {verbatim}, highlighted ones
% (```bash) as {Highlighting}; both are re-declared as breakable fancyvrb
% environments.  The continuation line is marked with a small hooked arrow.
\usepackage{fvextra}
\RecustomVerbatimEnvironment{verbatim}{Verbatim}{breaklines,breakanywhere,fontsize=\small}
\makeatletter
\@ifundefined{Highlighting}{}{%
  \RecustomVerbatimEnvironment{Highlighting}{Verbatim}{breaklines,breakanywhere,fontsize=\small,commandchars=\\\{\}}}
\makeatother
% inline code (pandoc writes it as \texttt{...}): never hyphenate -- an inserted
% hyphen could be read as part of a file name -- but allow a line break after
% an underscore, slash, comma, colon or equals sign, so that long paths and
% column lists wrap inside the text width and inside table columns.
% Inside a formula (e.g. the name of a notebook symbol in an equation) the same
% command must first leave math mode, which \text of amsmath does.
\ExplSyntaxOn
\tl_new:N \l__spt_code_tl
\cs_new_protected:Npn \__spt_texttt:n #1
  {
    \tl_set:Nn \l__spt_code_tl {#1}
    \regex_replace_all:nnN { ([\/\,\:\=]) } { \1 \c{allowbreak} } \l__spt_code_tl
    \regex_replace_all:nnN { \c{_} } { \c{_} \c{allowbreak} } \l__spt_code_tl
    \group_begin: \ttfamily \hyphenchar \font = -1 \scan_stop: \l__spt_code_tl \group_end:
  }
\RenewDocumentCommand \texttt { m }
  {
    \mode_if_math:TF { \text { \__spt_texttt:n {#1} } } { \__spt_texttt:n {#1} }
  }
\ExplSyntaxOff
% A figure is written in the Markdown source as an image paragraph followed by
% a caption paragraph ("**Figure S1.** ..."), so that the caption is also
% visible on GitHub.  An image that starts a paragraph is centered.
\usepackage{graphicx}% pandoc loads it only for documents that contain an image
\NewCommandCopy\sptincludegraphics\includegraphics
\RenewDocumentCommand\includegraphics{O{} m}{%
  \ifvmode\begin{center}\sptincludegraphics[#1]{#2}\end{center}%
  \else\sptincludegraphics[#1]{#2}\fi}
% keep figures where they are referenced
\usepackage{float}
\floatplacement{figure}{H}
% running header with the document title
\usepackage{fancyhdr}
\pagestyle{fancy}
\fancyhf{}
\fancyhead[L]{\small PROJECT spacetimes, Revision\_HS\_CJT\_and\_Kohn-Sham}
\fancyhead[R]{\small\thepage}
\renewcommand{\headrulewidth}{0.4pt}
% 16x16 matrices: tighter columns so that a full matrix fits the text width
\setlength{\arraycolsep}{3.2pt}
% The documents of this project contain long tagged symbols, file paths and
% identifiers that cannot be hyphenated.  A paragraph is therefore allowed to
% have looser lines before TeX gives up and lets a line stick out of the margin.
\AtBeginDocument{\tolerance=2000 \emergencystretch=5em \hbadness=2000}
\ifLuaTeX
  \usepackage{selnolig}  % disable illegal ligatures
\fi
\IfFileExists{bookmark.sty}{\usepackage{bookmark}}{\usepackage{hyperref}}
\IfFileExists{xurl.sty}{\usepackage{xurl}}{} % add URL line breaks if available
\urlstyle{same}
\hypersetup{
  pdftitle={The provenance of the Mathematica notebook superfluid.nb},
  pdfauthor={Prepared for Patrick Lee Nash, Ph.D.},
  colorlinks=true,
  linkcolor={blue},
  filecolor={Maroon},
  citecolor={blue},
  urlcolor={blue},
  pdfcreator={LaTeX via pandoc}}

\title{The provenance of the Mathematica notebook superfluid.nb}
\usepackage{etoolbox}
\makeatletter
\providecommand{\subtitle}[1]{% add subtitle to \maketitle
  \apptocmd{\@title}{\par {\large #1 \par}}{}{}
}
\makeatother
\subtitle{Provenance document \texttt{P3} of PROJECT spacetimes, folder
Revision\_HS\_CJT\_and\_Kohn-Sham}
\author{Prepared for Patrick Lee Nash, Ph.D.}
\date{November 21, 2026}

\begin{document}
\maketitle

\hypertarget{what-this-document-is}{%
\section{1. What this document is}\label{what-this-document-is}}

This document records how the Mathematica notebook
\texttt{superfluid.nb} of the folder
\texttt{Revision\_HS\_CJT\_and\_Kohn-Sham} was written, built and
evaluated, what it computes, which checks it makes and how each check is
repeated by a program in Rust and a program in Python, and which errors
were found and corrected while it was written. Every statement about the
results is taken from the log of the evaluation of the notebook in the
front end of Wolfram (the file \texttt{data/hck\_notebook\_log.txt}) and
from the reports of the Rust and the Python programs, by the program
\texttt{code/scripts/write\_notebook\_fragments.py}; no number of the
tables below is typed by hand.

\hypertarget{the-task-and-the-state}{%
\section{2. The task and the state}\label{the-task-and-the-state}}

\textbf{The task} (item 5 of the main task, in my words). A new,
stand-alone Mathematica notebook \texttt{superfluid.nb}, written in the
Wolfram language only, that evaluates without errors or warnings and is
an encyclopedia of the whole solution: every symbol, documented by
comments; every equation that can be executed; every solution; the
explanations, the figures, the graphs and the videos, each documented.
For every executable part of the notebook, executable pure Rust code and
Python code are part of the solution, with complete instructions, so
that users without a Wolfram license can follow; every time integration
in the notebook uses the method IDA inside the function \texttt{NDSolve}
(never a simple fixed-step integrator such as the fourth-order
Runge-Kutta method).

\textbf{The state of this version.} The notebook is written chapter by
chapter, and each chapter is committed as soon as it is complete. This
version contains the introduction, the register of the checks (chapter
0), Part I (items 1 to 5 of the problem: the conventions, the manifold
M31 and the Lorentz super-loop, the spinor field and theorem D1, the
gauge group GEE and its covariant derivatives, the Lagrangian and the
field equations, and the four analytic solutions), and the first section
of Part III (the Landau form of the effective potentials: its definition
with examples, its accuracy, the cause of its error and why the error
matters; the answer to the author's questions of 21 November 2026, which
he asked to be kept in the notebook), with 42 checks. The rest of Part
III, and Parts II and IV (the three methods and the superfluid, the
effective Lagrangians and the Mexican hat, the masses of the gauge
bosons, the surviving symmetry and the dynamics with IDA) follow in the
next versions. Until Part II is in the notebook, the section of Part III
takes the five computed polar curves of the effective potential from the
Rust program \texttt{sf\_part3} (a closed cell of data with the SHA-256
checksums of its files) and recomputes every number from them. Neither
part has a time integration.

\hypertarget{the-files}{%
\section{3. The files}\label{the-files}}

The files of the notebook, each with its SHA-256 checksum:

\begin{itemize}
\tightlist
\item
  \texttt{superfluid.nb}: the notebook, evaluated in the front end of
  Wolfram (inputs, printed checks, the table of all checks);\\
  \texttt{3b473f6015bb678f3cfa09f683c3c1cdc869fec08fbeacb184f8d3d2e837acf3}
\item
  \texttt{superfluid\_notebook.pdf}: the evaluated notebook printed by
  the front end, for readers without Wolfram;\\
  \texttt{15a48d963143b1ea34cc135adbb869bdc5d67cf5e4cbbb26b0591d1319651eea}
\item
  \texttt{wolfram/source/superfluid.source.wl}: the source of the
  notebook: the title, the introduction, chapter 0 (the register of the
  checks) and the final list of all checks; it includes the sources of
  the parts;\\
  \texttt{8446e7fb777d077765830ccbef7a42ac75ac3699c26c900840e6ed4dc41e4557}
\item
  \texttt{wolfram/source/part1.source.wl}: the source of Part I;\\
  \texttt{10bea1f4d47c952b00cf50af5183f2f3da3ee8af325ec485191108aaa789beb1}
\item
  \texttt{wolfram/source/part3\_landau.source.wl}: the source of the
  first section of Part III (the Landau form);\\
  \texttt{2fbe3ad5f084d3b4e038074f473f70bf5f3df715f2cadf9e48f930891a77d01e}
\item
  \texttt{wolfram/source/data\_part3.wl}: the values of Part III of the
  Rust program and its five polar curves, included in a closed cell of
  the notebook (written by \texttt{wolfram/tools/make\_data\_cells.py};
  its first lines name the data files and their SHA-256 checksums);\\
  \texttt{44deff1c0624be586db8a6899cf94e951150d20de5f3714a280ccf4f90d964e0}
\item
  \texttt{wolfram/tools/make\_data\_cells.py}: writes the data cells of
  the notebook from the data files of the Rust programs;\\
  \texttt{b1b63f10f2d8a206ac051e6b342c2c155ba8e1f55d589279945ab08515f477a6}
\item
  \texttt{wolfram/tools/build\_notebook.wls}: the builder: turns the
  source into the notebook (text cells with formulas, input cells parsed
  by the front end);\\
  \texttt{dd5e855c6fddcb0a00a6f3e2adecb5ced86c4cf52a8a48d4671eaa23fc63e1b0}
\item
  \texttt{wolfram/tools/sf\_formula.wl}: the reader of formulas: turns
  the TeX of the formulas of the texts into two-dimensional boxes;\\
  \texttt{a51dc3cc667c8b39366b517ba91b45502dca35e7eef95e9f796693c2853d6c12}
\item
  \texttt{wolfram/tools/guirun.wl}: the evaluation of the notebook in
  the front end, in a fresh kernel, with the log;\\
  \texttt{25f9b36cd6065a90db62a5db999da69a296e93f4ac3a42a878795ebafc9f2160}
\item
  \texttt{wolfram/tools/gui\_start.wl}: one of the three calls of the
  evaluation (Section 4);\\
  \texttt{210e94d1187b36befa5b731a81a643fc01f3b15b5092f34351969bf93d96762a}
\item
  \texttt{wolfram/tools/gui\_go.wl}: one of the three calls of the
  evaluation (Section 4);\\
  \texttt{1ad800de5df7a97cc38ed4755ce8f3769bce4c59adf2ed62bdb6c3b09a4e5a77}
\item
  \texttt{wolfram/tools/gui\_status.wl}: one of the three calls of the
  evaluation (Section 4);\\
  \texttt{0bc716d56ccc71443630e2d77ba6bfb00430660c80db499503a29bbde046967c}
\item
  \texttt{wolfram/tools/make\_dryrun.py}: writes the code of all input
  cells into one file for a test run by \texttt{wolframscript} without
  the front end;\\
  \texttt{f3ac70882c2c356d331bc5808683b7b56e535e9694ee39868c2b7cf3c2357f2f}
\item
  \texttt{data/hck\_notebook\_log.txt}: the log of the evaluation: the
  version of Wolfram, the counts of the cells, every message and printed
  line, one line per check;\\
  \texttt{4c26766542fc4ed82b5e165254db333c68616100cb5852cfd02704937ccb6a3e}
\item
  \texttt{data/hck\_notebook\_results.wl}: the register of the checks,
  written by the kernel at the end of the evaluation;\\
  \texttt{29b7817f4ed2488ffeeeef16eac62be393c9d10f36d6e8b334f769a7aa5ccf2b}
\item
  \texttt{code/scripts/write\_notebook\_fragments.py}: writes the tables
  and numbers of this document from the log and the reports;\\
  \texttt{e591d5ebc2fa730c7d621a73e0f8cbd95951bc26bec3926cf37a8ec585d5dc9a}
\end{itemize}

The builder, the reader of formulas and the evaluation script were first
written for the notebook \texttt{3+1super-loop.nb} of the project; they
are copied into this folder with their contexts renamed (from
\texttt{SL31} to \texttt{SFNB}), so that this folder holds every program
that makes its notebook. The reader of formulas there took its table of
the characters of TeX commands from a second package (the reader of
symbols of that notebook), which this notebook does not use; the table
is now part of \texttt{sf\_formula.wl} itself (error E2 of Section 9).

\hypertarget{how-the-notebook-is-built-and-evaluated}{%
\section{4. How the notebook is built and
evaluated}\label{how-the-notebook-is-built-and-evaluated}}

\textbf{For a reader with Wolfram.} Open \texttt{superfluid.nb} in
Wolfram (version 15 or later; this version was evaluated with 15.0.1 for
Mac OS X ARM (64-bit) (July 2, 2026)) and choose Evaluation
\textgreater{} Evaluate Notebook. Nothing outside the notebook is read.

\textbf{How the notebook is made from its source.} The source is a text
in the package format of the Wolfram language, in which every cell
begins with a marker line such as \texttt{(*\ ::Text::\ *)} or
\texttt{(*\ ::Input::\ *)}. On the author's computer:

\begin{verbatim}
wolframscript -file wolfram/tools/build_notebook.wls wolfram/source/superfluid.source.wl in/superfluid.nb \
    wolfram/tools/sf_formula.wl
\end{verbatim}

The builder replaces the lines
\texttt{(*\ @@INCLUDECELLS\ part1.source.wl@@\ *)} and
\texttt{(*\ @@INCLUDECELLS\ part3\_landau.source.wl@@\ *)} of the main
source by the cells of these parts, a line
\texttt{(*\ @@INCLUDE\ data\_part3.wl@@\ *)} inside an input cell by
that file, turns every formula of a text (TeX between dollar signs) into
boxes with the reader of formulas, lets the front end parse every input
cell (so that comments, line breaks and indentation stay as written),
and reports every formula that the reader cannot read as a warning. This
version: 108 cells (one of them a closed cell of data), 357 formulas in
the texts, 0 warnings.

\textbf{How it is evaluated.} The notebook is evaluated in the front end
of the running Wolfram application, exactly as by the menu command
Evaluate Notebook, but in a fresh kernel named \texttt{SFNBRun}, in
three calls (each returns at once):

\begin{enumerate}
\def\labelenumi{\arabic{enumi}.}
\tightlist
\item
  \texttt{gui\_start.wl}: opens the built notebook, sets its kernel, and
  starts the kernel with a hidden cell that makes the kernel append the
  number of every input line and the time to a progress file;
\item
  \texttt{gui\_go.wl}: evaluates all cells, and after them a hidden cell
  that writes the register of the checks
  (\texttt{SFNB}Test\texttt{\$Results}) to the results file and quits
  the kernel;
\item
  \texttt{gui\_status.wl}: saves the evaluated notebook, prints it to
  PDF, and writes the log: the version, the counts of the cells, every
  message and printed line, and one line per check.
\end{enumerate}

On the author's computer these calls are made by a small program that
runs inside the Wolfram application and executes requests left as files
in a work folder (\texttt{.git-transfer/mma/} of the author's copy of
the repository, which is not part of the repository); its handler
accepts a request \texttt{wls\ \textless{}script\textgreater{}} (run by
\texttt{wolframscript} in a fresh kernel) or
\texttt{gui\ \textless{}file\textgreater{}} (evaluated in the kernel of
the application). The first evaluation of this version made the three
calls with the same two Wolfram lines that the files
\texttt{gui\_start.wl}, \texttt{gui\_go.wl} and \texttt{gui\_status.wl}
now hold with comments.

\textbf{The dry run.} Before every build,
\texttt{python3\ wolfram/tools/make\_dryrun.py\ OUT.wls\ wolfram/source/superfluid.source.wl}
writes the code of all input cells, in order, into one file, and
\texttt{wolframscript\ -file\ OUT.wls} evaluates it in a plain kernel
(Part I and the section of Part III: 14 to 17 seconds); the only
difference from the notebook is that it prints the computing time after
every check. The dry run finds errors of the code in seconds; the
evaluation in the front end stays the only evidence for the results
reported here.

\hypertarget{the-evaluation-of-this-version}{%
\section{5. The evaluation of this
version}\label{the-evaluation-of-this-version}}

15.0.1 for Mac OS X ARM (64-bit) (July 2, 2026); 23 seconds; 20 input
cells, 20 of them evaluated; 5 output cells; 51 printed lines; \textbf{0
messages} (no error and no warning of the Wolfram language); verdict:
\textbf{ALL PASS (42 of 42), 0 messages}.

\hypertarget{the-checks-and-their-counterparts}{%
\section{6. The checks and their
counterparts}\label{the-checks-and-their-counterparts}}

Every check of the notebook carries the identifier of the same check in
the Rust program with the prefix W, and names its counterparts in the
Rust program (\texttt{rust/sfpart1}, program \texttt{sf\_part1}, for
Part I; \texttt{rust/sfpart3}, program \texttt{sf\_part3}, for Part III)
and in the Python program (\texttt{python/sf\_part1.py},
\texttt{python/sf\_part3.py}). The table gives, for each check, the
deviation measured by the notebook and by each counterpart (as their
reports state it; ``--'' for a check that is decided by facts, such as a
count or an exact identity). The program that writes the table stops
with an error if a counterpart is missing from its report or did not
pass.

\begin{longtable}[]{@{}
  >{\raggedright\arraybackslash}p{(\columnwidth - 12\tabcolsep) * \real{0.0825}}
  >{\raggedright\arraybackslash}p{(\columnwidth - 12\tabcolsep) * \real{0.0619}}
  >{\raggedright\arraybackslash}p{(\columnwidth - 12\tabcolsep) * \real{0.1134}}
  >{\raggedright\arraybackslash}p{(\columnwidth - 12\tabcolsep) * \real{0.1031}}
  >{\raggedright\arraybackslash}p{(\columnwidth - 12\tabcolsep) * \real{0.1340}}
  >{\raggedright\arraybackslash}p{(\columnwidth - 12\tabcolsep) * \real{0.1340}}
  >{\raggedright\arraybackslash}p{(\columnwidth - 12\tabcolsep) * \real{0.3711}}@{}}
\toprule\noalign{}
\begin{minipage}[b]{\linewidth}\raggedright
check
\end{minipage} & \begin{minipage}[b]{\linewidth}\raggedright
result
\end{minipage} & \begin{minipage}[b]{\linewidth}\raggedright
deviation in the notebook
\end{minipage} & \begin{minipage}[b]{\linewidth}\raggedright
tolerance
\end{minipage} & \begin{minipage}[b]{\linewidth}\raggedright
Rust: check, deviation
\end{minipage} & \begin{minipage}[b]{\linewidth}\raggedright
Python: check, deviation
\end{minipage} & \begin{minipage}[b]{\linewidth}\raggedright
statement
\end{minipage} \\
\midrule\noalign{}
\endhead
\bottomrule\noalign{}
\endlastfoot
WSM-1 & PASS & -- & -- & SM-1 (PASS), \(4.163\times10^{-16}\) & SM-1
(PASS), \(3.863\times10^{-16}\) & Standard Model with +i g: the
eigenvalues of the mass matrix are 0, M\_W\^{}2 (twice), M\_Z\^{}2, and
the photon sin(thetaW) W\^{}3 + cos(thetaW) B is massless {[}exact, with
symbolic g, g', v{]} \\
WA-2 & PASS & \(0\) & \(0\) & A-2 (PASS), \(0\) & A-2 (PASS), \(0\) &
bi-octonion units: e\_p\^{}2 = -1, e\_p e\_q = -e\_q e\_p, e\_j e\_4 =
e\_(4+j) (exact) \\
WA-3 & PASS & \(0\) & \(0\) & A-3 (PASS), \(0\) & A-3 (PASS), \(0\) &
N(xy) = N(x) N(y), 200 random Gaussian-integer bi-octonions (exact) \\
WA-4 & PASS & \(0\) & \(0\) & A-4 (PASS), \(0\) & A-4 (PASS), \(0\) & x
xbar = N(x) e\_0 (exact) \\
WA-5 & PASS & \(0\) & \(0\) & A-5 (PASS), \(0\) & A-5 (PASS), \(0\) &
alternativity (xx)y = x(xy), (yx)x = y(xx) (exact) \\
WA-6 & PASS & \(0\) & \(0\) & A-6 (PASS), \(0\) & A-6 (PASS), \(0\) &
the three Moufang identities (exact) \\
WA-7 & PASS & -- & -- & A-7 (PASS), -- & A-7 (PASS), -- & 168 of the 343
triples of imaginary units have a nonzero associator; (e1 e2) e4 - e1
(e2 e4) = 2 e7 {[}168 nonzero{]} \\
WA-8 & PASS & \(0\) & \(0\) & A-8 (PASS), \(1.776\times10^{-15}\) & A-8
(PASS), \(7.589\times10^{-15}\) & M(bb') = M(b)M(b'), det M = N,
M(b\^{}dag) = M(b)\^{}dag, inverse map (exact) \\
WA-10b & PASS & \(4.885\times10^{-15}\) & \(1.000\times10^{-10}\) &
A-10b (PASS), \(5.551\times10^{-15}\) & A-10b (PASS),
\(5.329\times10^{-15}\) & L\^{}T eta L = eta and det L = 1 (200 random
unit g) \\
WA-10c & PASS & -- & -- & A-10c (PASS), -- & A-10c (PASS), -- & L\_00
\textgreater= 1: orthochronous {[}200 random unit g{]} \\
WA-10d & PASS & \(5.329\times10^{-15}\) & \(1.000\times10^{-10}\) &
A-10d (PASS), \(4.441\times10^{-15}\) & A-10d (PASS),
\(7.105\times10^{-15}\) & L(gh) = L(g) L(h) and L(-g) = L(g) (200 random
pairs) \\
WA-12 & PASS & \(0\) & \(0\) & A-12 (PASS), \(0\) & A-12 (PASS), \(0\) &
spinor charges under the super-loop multiplications (exact) \\
WF-2 & PASS & \(0\) & \(0\) & F-2 (PASS), \(0\) & F-2 (PASS), \(0\) &
Clifford algebra, chirality, charge conjugation (exact) \\
WF-3a & PASS & \(1.777\times10^{-14}\) & \(1.000\times10^{-10}\) & F-3a
(PASS), \(1.599\times10^{-14}\) & F-3a (PASS), \(1.604\times10^{-14}\) &
S\^{}-1 gamma\^{}p S = sum\_q L\_pq gamma\^{}q (200 random g) \\
WF-3b & PASS & \(8.951\times10^{-16}\) & \(1.000\times10^{-10}\) & F-3b
(PASS), \(6.280\times10^{-16}\) & F-3b (PASS), \(7.448\times10^{-16}\) &
S\^{}dag gamma\^{}0 S = gamma\^{}0 (200 random g) \\
WF-3d & PASS & \(3.762\times10^{-15}\) & \(1.000\times10^{-10}\) & F-3d
(PASS), \(3.625\times10^{-15}\) & F-3d (PASS), \(3.827\times10^{-15}\) &
the Dirac matrices of the transformed frame are S gamma\^{}p S\^{}-1
(200 random g) \\
WF-5a & PASS & \(0\) & \(0\) & F-5a (PASS), \(0\) & F-5a (PASS), \(0\) &
D1: the Pauli-Guersey generators commute with the Lorentz generators
(exact) \\
WF-5c & PASS & -- & -- & F-5c (PASS), -- & F-5c (PASS), -- & D1: the
Lorentz-invariant bilinears of one point form a space of dimension 6;
those invariant also under the Pauli-Guersey SU(2), dimension 0
{[}dimensions 6, 0 (exact ranks){]} \\
WC-1 & PASS & \(5.551\times10^{-16}\) & \(1.000\times10^{-9}\) & C-1
(PASS), \(5.551\times10^{-16}\) & C-1 (PASS), \(3.736\times10^{-14}\) &
SU(2)\_v x U(1)\^{}F covariance of nabla on Psi (symbolic derivatives, 4
points) \\
WC-3 & PASS & \(1.178\times10^{-16}\) & \(1.000\times10^{-8}\) & C-3
(PASS), \(2.238\times10^{-16}\) & C-3 (PASS), \(2.755\times10^{-13}\) &
field strength W'\_pq = U W\_pq U\^{}dag (symbolic second derivatives, 2
points) \\
WC-4 & PASS & \(5.881\times10^{-16}\) & \(1.000\times10^{-9}\) & C-4
(PASS), \(9.037\times10^{-16}\) & C-4 (PASS), \(6.969\times10^{-14}\) &
isotriplet: nabla Phi = dPhi - g W x Phi + 2 i g\_F B Phi is covariant
(symbolic derivatives) \\
WC-8 & PASS & \(1.256\times10^{-15}\) & \(1.000\times10^{-9}\) & C-8
(PASS), \(2.809\times10^{-15}\) & C-8 (PASS), \(5.043\times10^{-14}\) &
local SL(2,C): gamma\^{}p{[}f'{]} (d\_p + Gamma'\_p)(S psi) = S
gamma\^{}p d\_p psi (symbolic derivatives) \\
WE-2 & PASS & \(0\) & \(0\) & E-2 (PASS), \(0\) & E-2 (PASS), \(0\) &
scalar contact term: {[}Psi\_e, H\_S{]} = G\^{}S (Psi-bar
Psi)(gamma\^{}0 Psi)\_e (exact, 256 states) \\
WE-3 & PASS & \(0\) & \(0\) & E-3 (PASS), \(3.553\times10^{-15}\) & E-3
(PASS), \(3.553\times10^{-15}\) & vector contact term (exact, 256
states) \\
WE-4 & PASS & \(0\) & \(0\) & E-4 (PASS), \(0\) & E-4 (PASS), \(0\) &
pairing term: {[}Psi\_e, H\_P{]} = -(G/2) sum (K\^{}dag Psi\^{}*)\_e Pi
(exact, 256 states) \\
WE-14 & PASS & -- & -- & E-14 (PASS), \(2.331\times10^{-15}\) & E-14
(PASS), -- & symbolic Euler-Lagrange equations of the Yang-Mills
Lagrangian with the source: d\_p W\^{}(a,p p') - g eps\_abc W\^{}b\_p
W\^{}(c,p p') = g j\^{}(a,p') (12 components) {[}12 of 12 equations
identical (EulerEquations; signs \{1, 1, 1, 1, 1, 1, 1, 1, 1, 1, 1,
1\}){]} \\
WAS1-3 & PASS & \(1.776\times10^{-15}\) & \(1.000\times10^{-12}\) &
AS1-3 (PASS), \(2.887\times10^{-15}\) & AS1-3 (PASS),
\(1.554\times10^{-15}\) & spin-isospin locked field: four analytic
branches (61 momenta) \\
WAS2-2 & PASS & \(4.828\times10^{-3}\) & \(1.000\times10^{-2}\) & AS2-2
(PASS), \(4.828\times10^{-3}\) & AS2-2 (PASS), \(4.828\times10^{-3}\) &
Landau levels, Wilson lattice spacing 0.025 (largest relative deviation
of the seven levels) \\
WAS2-4 & PASS & \(6.450\times10^{-5}\) & \(1.000\times10^{-3}\) & AS2-4
(PASS), \(6.450\times10^{-5}\) & AS2-4 (PASS), \(6.450\times10^{-5}\) &
Landau levels, Richardson extrapolation 2 E(0.025) - E(0.05) \\
WAS3-1 & PASS & \(2.665\times10^{-15}\) & \(1.000\times10^{-12}\) &
AS3-1 (PASS), \(3.775\times10^{-15}\) & AS3-1 (PASS),
\(2.665\times10^{-15}\) & BdG spectrum of unitary pair fields (50
random) \\
WAS3-4 & PASS & \(3.005\times10^{-11}\) & \(1.000\times10^{-9}\) & AS3-4
(PASS), -- & AS3-4 (PASS), \(3.005\times10^{-11}\) & the gap at G = 22.5
agrees with the Rust value 0.2094933033 \\
WAS4-5a & PASS & -- & -- & AS4-5a (PASS), -- & AS4-5a (PASS), -- & three
bound states below the continuum for each spacing {[}counts \{3, 3,
3\}{]} \\
WAS4-5b & PASS & \(1.116\times10^{-5}\) & \(1.000\times10^{-3}\) &
AS4-5b (PASS), \(1.116\times10^{-5}\) & AS4-5b (PASS),
\(1.116\times10^{-5}\) & kink bound states, Richardson extrapolation \\
WAS4-3 & PASS & -- & -- & AS4-3 (PASS), \(2.776\times10^{-17}\) & AS4-3
(PASS), -- & symbolic: (d/dx + Delta) sech\^{}2(x/ell) = 0 and (-d\^{}2
+ Delta\^{}2 - Delta') sech tanh = (3/16) sech tanh {[}residuals 0,
0{]} \\
WP3-11 HS\_1\_1 & PASS & \(2.492\times10^{-11}\) &
\(1.000\times10^{-7}\) & P3-11 HS\_1\_1 (PASS), \(1.103\times10^{-14}\)
& YP3-11 HS\_1\_1 (PASS), \(9.618\times10^{-12}\) & HS (1, 1): the
accuracy of the Landau form recomputed by LeastSquares (the
coefficients, the minimum, depth and curvature of the polynomial, the
residuals of the fits with Phi\^{}2 .. Phi\^{}(2K), K = 2 .. 6, the fits
on three intervals, the sign changes) = Rust (largest relative
difference) \\
WP3-11 HS\_0\_1 & PASS & \(4.409\times10^{-11}\) &
\(1.000\times10^{-7}\) & P3-11 HS\_0\_1 (PASS), \(5.616\times10^{-13}\)
& YP3-11 HS\_0\_1 (PASS), \(1.478\times10^{-11}\) & HS (0, 1): the
accuracy of the Landau form recomputed by LeastSquares (the
coefficients, the minimum, depth and curvature of the polynomial, the
residuals of the fits with Phi\^{}2 .. Phi\^{}(2K), K = 2 .. 6, the fits
on three intervals, the sign changes) = Rust (largest relative
difference) \\
WP3-11 HS\_0\_0 & PASS & \(4.051\times10^{-11}\) &
\(1.000\times10^{-7}\) & P3-11 HS\_0\_0 (PASS), \(8.612\times10^{-13}\)
& YP3-11 HS\_0\_0 (PASS), \(5.701\times10^{-11}\) & HS (0, 0): the
accuracy of the Landau form recomputed by LeastSquares (the
coefficients, the minimum, depth and curvature of the polynomial, the
residuals of the fits with Phi\^{}2 .. Phi\^{}(2K), K = 2 .. 6, the fits
on three intervals, the sign changes) = Rust (largest relative
difference) \\
WP3-11 KS & PASS & \(3.482\times10^{-9}\) & \(1.000\times10^{-7}\) &
P3-11 KS (PASS), \(7.503\times10^{-13}\) & YP3-11 KS (PASS),
\(3.104\times10^{-9}\) & KS: the accuracy of the Landau form recomputed
by LeastSquares (the coefficients, the minimum, depth and curvature of
the polynomial, the residuals of the fits with Phi\^{}2 .. Phi\^{}(2K),
K = 2 .. 6, the fits on three intervals, the sign changes) = Rust
(largest relative difference) \\
WP3-11 CJT & PASS & \(6.816\times10^{-9}\) & \(1.000\times10^{-7}\) &
P3-11 CJT (PASS), \(2.943\times10^{-13}\) & YP3-11 CJT (PASS),
\(1.216\times10^{-9}\) & CJT: the accuracy of the Landau form recomputed
by LeastSquares (the coefficients, the minimum, depth and curvature of
the polynomial, the residuals of the fits with Phi\^{}2 .. Phi\^{}(2K),
K = 2 .. 6, the fits on three intervals, the sign changes) = Rust
(largest relative difference) \\
WP3-8 & PASS & \(1.554\times10^{-15}\) & \(1.000\times10^{-13}\) & P3-8
(PASS), \(2.442\times10^{-15}\) & YP3-8 (PASS), \(1.332\times10^{-15}\)
& the pair field enters the quasi-particle energies like a mass: the BdG
matrix with N = (M - m) beta (x) 1\_2 and P = Phi\_a (K\^{}3)\^{}dag has
the eigenvalues -+(q\^{}2 + M\^{}2 + Phi\_a\^{}2)\^{}(1/2), eight times
each (20 random q, M, Phi\_a) \\
WP3-10 HS\_1\_1 & PASS & \(2.646\times10^{-7}\) & \(1.000\times10^{-6}\)
& P3-10 HS\_1\_1 (PASS), \(8.786\times10^{-9}\) & YP3-10 HS\_1\_1
(PASS), \(3.072\times10^{-11}\) & HS (1, 1): the coefficient of Phi\^{}2
of the Taylor series at the origin from the first six points of the
curve = the Rust value a\^{}Tay = -0.00616243 (relative) \\
WP3-10 CJT & PASS & \(2.377\times10^{-8}\) & \(1.000\times10^{-6}\) &
P3-10 CJT (PASS), \(7.425\times10^{-8}\) & YP3-10 CJT (PASS),
\(2.872\times10^{-12}\) & CJT: the coefficient of Phi\^{}2 of the Taylor
series at the origin from the first six points of the curve = the Rust
value a\^{}Tay = -0.00429666 (relative) \\
\end{longtable}

\textbf{Table 1.} The checks of the notebook. Data:
\texttt{data/hck\_notebook\_log.txt},
\texttt{data/hck\_part1\_checks.txt},
\texttt{data/hck\_part1\_checks\_python.txt},
\texttt{data/hck\_part3\_checks.txt},
\texttt{data/hck\_part3\_checks\_python.txt}.

\hypertarget{the-methods-of-the-notebook-and-of-its-counterparts}{%
\section{7. The methods of the notebook and of its
counterparts}\label{the-methods-of-the-notebook-and-of-its-counterparts}}

The three implementations compute the same quantities by different
means, so that an error of one method does not repeat itself in the
others:

\begin{itemize}
\tightlist
\item
  \textbf{bi-quaternions and bi-octonions} (WA-2 to WA-12): the notebook
  uses Gaussian-integer coefficients and exact arithmetic, so every
  identity holds with the deviation exactly 0; the Rust and the Python
  programs use small integer coefficients in double precision (also
  exact) and random floating-point values where the notebook uses exact
  ones;
\item
  \textbf{Lorentz and spinor matrices} (WA-10, WF-3): floating point in
  all three, with different random group elements (the seeds of the
  three generators differ);
\item
  \textbf{theorem D1} (WF-5): exact ranks of exact matrices in the
  notebook; ranks by the singular values in Python, and by the Rust
  program's own elimination;
\item
  \textbf{covariant derivatives} (WC-1 to WC-8): symbolic derivatives
  (the function \texttt{D}) of the random smooth fields in the notebook;
  automatic differentiation in Rust; central differences of eighth order
  in Python;
\item
  \textbf{the field equations as commutators} (WE-2 to WE-4): exact
  operator identities on the 256 states of one point in the notebook
  (exact couplings \(-8\), \(4/5\), \(45/2\), deviation exactly 0);
  double precision in the two programs;
\item
  \textbf{the Yang-Mills equations} (WE-14): the function
  \texttt{EulerEquations} of the package \texttt{VariationalMethods} in
  the notebook; the program sympy in Python; in Rust, exact derivatives
  of the Lagrangian density, a polynomial of degree four in the fields
  and their first derivatives (five-point differences in its arguments,
  exact for such a polynomial, and automatic differentiation along the
  coordinates);
\item
  \textbf{the conventions of the Standard Model} (WSM-1): exact, with
  symbolic couplings, in the notebook (the characteristic polynomial and
  the null vector of the mass matrix); the four eigenvectors with 20
  values of the couplings in Rust; the eigenvalues by LAPACK with 20
  random values in Python;
\item
  \textbf{the analytic solutions} (WAS1-3 to WAS4-3): dense eigenvalues
  of the Hermitian matrices by the function \texttt{Eigenvalues} (up to
  4800 rows), the gap equation by \texttt{FindRoot} with 30 digits, and
  the profiles of the kink by \texttt{FullSimplify} in the notebook; the
  Rust program's own eigensolvers (inertia counting and bisection of
  block-tridiagonal matrices) and the band solver of LAPACK in Python;
\item
  \textbf{the accuracy of the Landau form} (WP3-8, WP3-10, WP3-11): the
  least squares by the function \texttt{LeastSquares} on scaled columns
  in the notebook, by Householder reflections in Rust, and by the
  singular-value decomposition of numpy in Python; the quasi-particle
  energies by \texttt{Eigenvalues} of the \(16\times16\) matrices (the
  Rust program's own eigensolver, LAPACK in Python); the Taylor
  coefficient at the origin from the first six points of the computed
  curve in the notebook, and from three states with small applied fields
  and the Richardson extrapolation in the two programs; the integral
  with the term \(\tfrac1{16}s^2\ln s\) exactly, by \texttt{Integrate},
  \texttt{FullSimplify} and \texttt{Series}, in the notebook alone (it
  is shown, not checked).
\end{itemize}

\hypertarget{the-symbols-of-the-code-of-part-i}{%
\section{8. The symbols of the code of Part
I}\label{the-symbols-of-the-code-of-part-i}}

Every name of the code that stands for an object of the physics is
introduced in the notebook with a comment that gives its printed form
and its meaning. The table lists them in the order of the notebook. (The
names \texttt{dA2}, \texttt{dF3a}, \texttt{dC1}, \texttt{dE2} and their
like hold the measured deviation of the check whose identifier they
contain; the names that end in \texttt{T}, such as \texttt{uT} or
\texttt{psiT}, are the random test fields of Section I.4.)

\begin{longtable}[]{@{}
  >{\raggedright\arraybackslash}p{(\columnwidth - 4\tabcolsep) * \real{0.3333}}
  >{\raggedright\arraybackslash}p{(\columnwidth - 4\tabcolsep) * \real{0.3333}}
  >{\raggedright\arraybackslash}p{(\columnwidth - 4\tabcolsep) * \real{0.3333}}@{}}
\toprule\noalign{}
\begin{minipage}[b]{\linewidth}\raggedright
name in the code
\end{minipage} & \begin{minipage}[b]{\linewidth}\raggedright
printed form
\end{minipage} & \begin{minipage}[b]{\linewidth}\raggedright
meaning
\end{minipage} \\
\midrule\noalign{}
\endhead
\bottomrule\noalign{}
\endlastfoot
\texttt{check}, \texttt{checkTrue}, \texttt{sci}, \texttt{checkHook} &
-- & the functions that record and print a check, the number format of
the reports, and the hook used by the dry run \\
\texttt{mLL} & \(m_{\rm LL}\) & the mass of the spinor field, \(1/5\) in
units of the cutoff \\
\texttt{sig{[}{[}p+1{]}{]}}, \texttt{sigbar{[}{[}p+1{]}{]}} &
\(\sigma^{\rm LL}_p\), \(\bar\sigma_p\) & the unit matrix and the Pauli
matrices; \(\sigma^{\rm LL}_0\) and \(-\sigma^{\rm LL}_j\) \\
\texttt{eLL}, \texttt{etaLL}, \texttt{levi} & \(\mathsf E_{\rm LL}\),
\(\eta^{\rm LL}_{pp'}\), \(\epsilon^{\rm LL}\) & \(i\sigma^{\rm LL}_2\),
the Minkowski metric, the Levi-Civita symbol \\
\texttt{gam{[}p{]}}, \texttt{gch}, \texttt{cLL}, \texttt{betaLL},
\texttt{alphaLL{[}j{]}} & \(\gamma^p_{\rm LL}\),
\(\gamma^{\rm ch}_{\rm LL}\), \(C_{\rm LL}\), \(\beta_{\rm LL}\),
\(\alpha^j_{\rm LL}\) & the Dirac matrices of the chiral basis,
chirality, charge conjugation, and the matrices of the Dirac
Hamiltonian \\
\texttt{d8}, \texttt{i8} & \(A\otimes\mathsf 1_2\),
\(\mathsf 1_4\otimes t\) & a Dirac matrix and an isospin matrix acting
on the eight components \\
\texttt{tLL{[}iota{]}}, \texttt{kLL{[}iota{]}}, \texttt{hLL{[}q,\ m{]}}
& \(T^\iota_{\rm LL}\), \(\mathsf K^\iota_{\rm LL}\),
\(\mathsf h_{\rm LL}(\vec q_{\rm LL})\) & the isospin generators, the
matrices of the scalar spin-zero pairs, the free Dirac Hamiltonian \\
\texttt{qMul}, \texttt{qBar}, \texttt{qDag}, \texttt{qNorm} &
\(\mathbf b\mathbf b'\), \(\bar{\mathbf b}\), \(\mathbf b^\dagger\),
\(N_{\rm LL}(\mathbf b)\) & the product, the conjugations and the
complex norm of bi-quaternions \\
\texttt{mLLmat}, \texttt{biqOf} & \(M_{\rm LL}(\mathbf b)\) & the matrix
of a bi-quaternion and the inverse map \\
\texttt{oMul}, \texttt{oUnit{[}p{]}}, \texttt{oBar}, \texttt{oNorm},
\texttt{oMax} & \(\mathbf x\mathbf y\), \(\mathbf e_p\),
\(\bar{\mathbf x}\), \(N_{\rm LL}(\mathbf x)\) & the product of Cayley
and Dickson, the units, the conjugation and the norm of bi-octonions,
and the size of a difference \\
\texttt{expTraceless}, \texttt{lorentzGenerator},
\texttt{lorentzMatrix}, \texttt{sSpin}, \texttt{sGen},
\texttt{randUnitBiq} & \(\exp\), \(M_{\rm LL}(\mathbf u)\),
\(\mathsf L_{\rm LL}(\mathbf g)\), \(\mathsf S_{\rm LL}(\mathbf g)\),
\(\mathsf s_{\rm LL}(\mathbf u)\) & the exponential of a traceless
\(2\times2\) matrix, a Lorentz generator, the Lorentz matrix, the spinor
matrix, the action of a generator on the Dirac index, a random unit
bi-quaternion \\
\texttt{fockOps}, \texttt{cOps}, \texttt{cdOps} & \(\Psi_{{\rm LL},c}\),
\(\Psi^\dagger_{{\rm LL},c}\) & the annihilation and creation operators
of the eight components of one point (256 states) \\
\texttt{gSM}, \texttt{gpSM}, \texttt{vSM}, \texttt{phiSM}, \texttt{xSM},
\texttt{mSM}, \texttt{photonSM} & \(g\), \(g'\), \(v\),
\(\langle\Phi^{\rm SM}\rangle\), -- , -- , \(A_\mu\) & the couplings and
the vacuum value of the Standard Model, the vacuum, the vectors \(X_a\)
of \(D_\mu\langle\Phi^{\rm SM}\rangle=\sum_aV^a_\mu X_a\), the mass
matrix of \((W^1,W^2,W^3,B)\), the photon \\
\texttt{lorGens}, \texttt{pgGens}, \texttt{amatD1}, \texttt{basisD1},
\texttt{invDim} & -- & theorem D1: the Lorentz and the Pauli-Gürsey
generators on \((\psi,\psi^\dagger)\), the map to the doublet of Weyl
spinors, the antisymmetric bilinears, the dimension of the invariant
ones \\
\texttt{xs}, \texttt{gvT}, \texttt{gfT} & \(x^p_{\rm LL}\),
\(g^{\rm v}_{\rm LL}\), \(g^{\rm F}_{\rm LL}\) & the coordinates as
symbols and the couplings of the test, \(0.83\) and \(0.47\) \\
\texttt{uT}, \texttt{thT}, \texttt{wT{[}p{]}}, \texttt{bT{[}p{]}},
\texttt{wTp{[}p{]}}, \texttt{bTp{[}p{]}} & \(\mathsf U_{\rm LL}\),
\(\theta^{\rm F}_{\rm LL}\), \(\mathsf W_{{\rm SF},p}\),
\(B_{{\rm SF},p}\), \(\mathsf W'_{{\rm SF},p}\), \(B'_{{\rm SF},p}\) &
the test fields of the covariance checks and their gauge transforms \\
\texttt{psiT}, \texttt{psiTp}, \texttt{covPsi}, \texttt{phiT},
\texttt{phiTp}, \texttt{rotT}, \texttt{covPhi}, \texttt{fieldStrength} &
\(\Psi_{\rm LL}\), \(\nabla^{\rm GEE}_p\Psi_{\rm LL}\), \(\vec\Phi\),
\(\mathsf R^{\rm v}_{\rm LL}\), \(\nabla_p\vec\Phi\),
\(\mathsf W_{{\rm SF},pp'}\) & the test spinor and its transform, the
covariant derivative, the isotriplet and its rotation, its covariant
derivative, the field strength \\
\texttt{mF}, \texttt{sF}, \texttt{sFinv}, \texttt{psi4} &
\(M_{\rm LL}(\mathbf g(x))\), \(\mathsf S_{\rm LL}(\mathbf g(x))\) & the
field of unit bi-quaternions of the local \(SL(2,\mathbb C)\) check, its
spinor matrix and inverse, and the test Dirac spinor \\
\texttt{gScal}, \texttt{gVec}, \texttt{gPair}, \texttt{hS}, \texttt{hV},
\texttt{hP}, \texttt{piOps} & \(G^{\rm S}_{\rm SF}\),
\(G^{\rm V}_{\rm SF}\), \(G_{\rm LL}\), -- , -- , -- ,
\(\Pi^\iota_{\rm LL}\) & the exact couplings \(-8\), \(4/5\), \(45/2\)
of the operator checks, the Hamiltonians of the three contact terms of
one point, the pair operators \\
\texttt{wF}, \texttt{jF}, \texttt{fYM}, \texttt{lYM}, \texttt{eqsYM},
\texttt{gYM} & \(W^\iota_{{\rm SF},p}\), \(j^{\iota,p}_{\rm SF}\),
\(W^\iota_{{\rm SF},pp'}\), -- , -- , \(g^{\rm v}_{\rm LL}\) & the
symbolic gauge fields and currents, the field strengths, the Lagrangian
density of the Yang-Mills check, its Euler-Lagrange equations, the
coupling \\
\texttt{hPW}, \texttt{i1}, \texttt{gapAS3} & -- , \(I_1\), \(\Delta\) &
the plane-wave matrix of AS1, the integral of the gap equation of AS3,
the gap at \(G=22.5\) \\
\texttt{lfVec}, \texttt{lfNorm2}, \texttt{lfIdentity}, \texttt{lfPolar}
& \(\vec\Phi\), \(\lvert\vec\Phi\rvert^2\) & Part III: a complex
3-vector with symbolic real and imaginary parts (\texttt{lfRe{[}k{]}},
\texttt{lfIm{[}k{]}}), its squared norm, the identity
\(\lvert\vec\Phi\rvert^4-\lvert\vec\Phi\cdot\vec\Phi\rvert^2=\lvert\vec\Phi\times\vec\Phi^*\rvert^2\)
and its value on a polar field (\texttt{lfTh}, \texttt{lfR},
\texttt{lfT1}, \texttt{lfT2}: its phase, modulus and direction) \\
\texttt{rustPart3}, \texttt{rustEffPot{[}tag{]}} & -- & Part III: the
values of the Rust program \texttt{sf\_part3} and its five computed
polar curves (applied field, field, energy per volume, Dirac mass) \\
\texttt{lfTags}, \texttt{lfName}, \texttt{lfPhi0}, \texttt{lfE0},
\texttt{lfCurve}, \texttt{lfFitted} & \(X\), \(\Phi^X_{0,{\rm SF}}\),
\(\mathcal V^X_{\rm SF}(\Phi^X_{0,{\rm SF}})\) & Part III: the five
methods, their names, the vacuum field, the energy of the minimum (minus
the depth), the polar curve \(\{\Phi,\mathcal V^X_{\rm SF}\}\), and its
points with \(\Phi\le\) a given multiple of \(\Phi^X_{0,{\rm SF}}\) \\
\texttt{evenFit}, \texttt{landauMin}, \texttt{lfResults},
\texttt{lfKeys}, \texttt{lfRows}, \texttt{pct} & \(\sum_kc_k\Phi^{2k}\),
\(\Phi_{\rm L}\) & Part III: the least-squares fit with the powers
\(\Phi^2,\dots,\Phi^{2K}\), the minimum of the Landau polynomial with
its value and second derivative, all numbers of the accuracy for each
method, the names compared with the Rust values, the rows of the printed
table, a number in per cent \\
\texttt{lfColors}, \texttt{lfShown}, \texttt{lfLegend} & -- & Part III:
the colors, the four distinct curves (KS and CJT coincide) and the
legend of the figure \\
\texttt{lfBdG}, \texttt{lfQpDev} &
\(\mathsf H_{\rm SF}(\vec q_{\rm LL})\) & Part III: the matrix of
Bogoliubov and de Gennes with a Dirac mass and an applied pair field,
and the deviation of its eigenvalues from \(\pm E\) \\
\texttt{lfq}, \texttt{lfs}, \texttt{lfLam}, \texttt{lfInt},
\texttt{lfClosed}, \texttt{lfSeries} & \(q\), \(s\),
\(\Lambda_{\rm LL}\) & Part III: the variables of the integral
\(\int_0^{\Lambda}q^2(q^2+s)^{1/2}dq\), its value, its closed form and
its expansion about \(s=0\) \\
\texttt{lfTaylor} & \(a^{X,{\rm Tay}}_{\rm SF}\) & Part III: the Taylor
coefficient of \(\lvert\vec\Phi\rvert^2\) at the origin, estimated from
the first six points of the curve \\
\end{longtable}

\hypertarget{errors-and-corrections-found-while-this-version-was-written}{%
\section{9. Errors and corrections found while this version was
written}\label{errors-and-corrections-found-while-this-version-was-written}}

\begin{enumerate}
\def\labelenumi{\arabic{enumi}.}
\item
  \emph{A comment that ended early (E1).} Two comments of the code of
  theorem D1 contained the characters \texttt{\^{}*)} (in
  \texttt{s(u)\^{}*} and \texttt{(psi\^{}+)\^{}*}), which end a comment
  of the Wolfram language; the first dry run stopped with a syntax
  error. The comments now write the complex conjugate as
  \texttt{conj(...)}.
\item
  \emph{The table of the characters of TeX commands (E2).} The reader of
  formulas took this table from the reader of symbols of the notebook
  \texttt{3+1super-loop.nb}; without it the reader would not know the
  Greek letters. The table is now part of \texttt{sf\_formula.wl}.
\item
  \emph{A name that resembled a forbidden word (E3).} The first
  evaluated notebook named the coupling of the pairing term with three
  letters that the scan of every commit rejects (the scan refuses a
  short list of names in the files of the project). The couplings are
  now named \texttt{gScal}, \texttt{gVec} and \texttt{gPair}, and the
  notebook was built and evaluated again.
\item
  \emph{Two checks with one identifier (E4).} The table of Section 6
  compares every check with its counterparts by their identifiers. It
  showed that the Python program had used the identifier E-10 for its
  symbolic derivation of the Yang-Mills equations, while E-10 of the
  Rust program is a different check (the self-energy of Hartree and Fock
  is the derivative of the energy). The Python check is now E-14, the
  Rust program has a new check E-14 of the same statement (exact
  derivatives of the polynomial Lagrangian density), and the notebook's
  check is WE-14. The manuscript (Section 5.2 and Section 17.5) states
  this.
\item
  \emph{A check without counterparts (E5).} The notebook's check of the
  conventions of the Standard Model had at first no counterpart in the
  two programs; the check SM-1 was added to both (Section 1 of the
  manuscript).
\item
  \emph{The printed time of the dry run (E6).} The first dry run printed
  the computing times as unevaluated number formats; the dry run now
  prints rounded numbers. This affects only the dry run.
\item
  \emph{A figure wider than the printed page (E7).} The first evaluation
  of the section of Part III placed the two panels of its figure side by
  side, 880 points wide; the printed copy cut the right panel. The
  panels now stand one above the other, each 460 points wide, and the
  notebook was built and evaluated again.
\end{enumerate}

\hypertarget{what-follows}{%
\section{10. What follows}\label{what-follows}}

The rest of Part III, and Parts II and IV of the notebook (the three
methods, the effective Lagrangians, the masses of the gauge bosons, the
surviving symmetry, the dynamics with IDA inside \texttt{NDSolve}, the
figures and the videos) are added in the same way: the checks of the
Python programs \texttt{sf\_part2.py}, \texttt{sf\_part3.py} and
\texttt{sf\_part4.py}, each with its counterparts, and this document is
extended with each part.

\end{document}
